\documentclass[12pt, a4paper]{article}

% 引入必要的宏包
\usepackage[UTF8]{ctex}     % 中文支持
\usepackage{amsmath, amssymb, amsthm} % 数学符号和环境
\usepackage{geometry}       % 页面设置
\usepackage{enumitem}       % 列表定制
\usepackage{fancyhdr}       % 页眉页脚
\usepackage{cases}          % 分段函数增强

% 页面边距设置
\geometry{left=2.5cm, right=2.5cm, top=2.5cm, bottom=2.5cm}

% 宏定义：方便排版选择题选项
% 横向排列的选项 (4列)
\newcommand{\choicesFour}[4]{
    \begin{enumerate}[label=\Alph*., itemsep=0pt, parsep=0pt, topsep=5pt, labelsep=0.5em]
        \begin{minipage}{0.22\linewidth} \item #1 \end{minipage}
        \begin{minipage}{0.22\linewidth} \item #2 \end{minipage}
        \begin{minipage}{0.22\linewidth} \item #3 \end{minipage}
        \begin{minipage}{0.22\linewidth} \item #4 \end{minipage}
    \end{enumerate}
}
% 横向排列的选项 (2列)
\newcommand{\choicesTwo}[4]{
    \begin{enumerate}[label=\Alph*., itemsep=0pt, parsep=0pt, topsep=5pt, labelsep=0.5em]
        \begin{minipage}{0.45\linewidth} \item #1 \end{minipage}
        \begin{minipage}{0.45\linewidth} \item #2 \end{minipage}
        \begin{minipage}{0.45\linewidth} \item #3 \end{minipage}
        \begin{minipage}{0.45\linewidth} \item #4 \end{minipage}
    \end{enumerate}
}
% 纵向排列的选项 (1列)
\newcommand{\choices}[4]{
    \begin{enumerate}[label=\Alph*., itemsep=0pt, parsep=0pt, topsep=5pt, labelsep=0.5em]
        \item #1
        \item #2
        \item #3
        \item #4
    \end{enumerate}
}

% 填空题下划线
\newcommand{\blank}[1]{\underline{\hspace{#1}}}

\title{\textbf{《高等数学》必刷习题——提高版}}
\author{}
\date{}

\begin{document}

\section*{第十二章 无穷级数提高十五题}

\begin{enumerate}
    \item 有下列关于数项级数的命题，其中正确的个数为 \blank{1cm}。
    \begin{enumerate}[label=(\arabic*)]
        \item 若 $ \lim_{n\to\infty}u_{n}\neq0 $ ,则 $ \sum_{n=1}^{\infty}u_{n} $ 必发散.
        \item 若 $ u_{n}\geq0,u_{n}\geq u_{n+1}(n=1,2,3\ldots) $ 且 $ \lim_{n\to\infty}u_{n}=0, $ 则 $ u_{n} $ 必收敛.
        \item 若 $ \sum_{n=1}^{\infty}u_{n} $ 收敛，则 $ \sum_{n=1}^{\infty}|u_{n}| $ 必收敛.
        \item 若 $ \sum_{n=1}^{\infty}u_{n} $ 收敛于 $s$,则任意改变该级数项的位置所得到的新的级数仍收敛于 $s$.
    \end{enumerate}
    \choicesFour{0}{1}{2}{3}

    \item 下列级数中收敛的级数是 \blank{1cm}。
    \choicesFour{$ \sum_{n=1}^{\infty}\frac{1}{3^{n}} $}{$ \sum_{n=1}^{\infty}\frac{1}{n+1} $}{$ \sum_{n=1}^{\infty}\frac{3^{n}}{2^{n}} $}{$ \sum_{n=1}^{\infty}\frac{1}{\sqrt{n+1}} $}

    \item 级数 $ \sum_{n=1}^{\infty} (-1)^{n} \frac{1}{n^{2}} $ 是 \blank{1cm}。
    \choicesFour{发散的}{条件收敛}{绝对收敛}{收敛性不能确定}

    \item 下列级数中为条件收敛的级数是 \blank{1cm}。
    \choicesTwo
    {$ \sum_{n=1}^{\infty}(-1)^{n}\frac{n}{n+1} $}
    {$ \sum_{n=1}^{\infty}(-1)^{n}\sqrt{n} $}
    {$ \sum_{n=1}^{\infty}(-1)^{n}\frac{1}{n^{2}} $}
    {$ \sum_{n=1}^{\infty}(-1)^{n}\frac{1}{\sqrt{n}} $}

    \item 幂级数 $ \sum_{n=0}^{\infty}\frac{x^{n}}{n!} $ 的收敛半径是 \blank{1cm}。
    \choicesFour{1}{0}{$ +\infty $}{$ \frac{1}{2} $}

    \item 求幂级数 $ \sum_{n=1}^{\infty} \frac{(x-5)^{n}}{\sqrt{n}} $ 的收敛域 \blank{2cm}。

    \item 判别级数的收敛性，级数是 $ \sum_{n=1}^{\infty}\frac{n^{n}}{(n!)^{2}} $ \blank{2cm}。(收敛还是发散)

    \item 判别级数的收敛性，级数是 $ \sum_{n=0}^{\infty} (-1)^{n} \frac{(n+1)!}{n^{n+1}} $ \blank{2cm}。(绝对收敛，条件收敛还是发散)

    \item 求幂级数 $ \sum_{n=0}^{\infty} 10^{n}(x-1)^{n} $ 的收敛域 \blank{2cm}。

    \item 求幂级数 $ \sum_{n=0}^{\infty} \frac{x^{n}}{2^{n-1}} $ 的和函数 \blank{2cm}。

    \item 求幂级数 $ \sum_{n=0}^{\infty}\frac{x^{n}}{\sqrt{n+1}} $ 的收敛域。

    \item 求幂级数 $ \sum_{n=0}^{\infty} \frac{x^{n+2}}{n+1} $ 的收敛域及和函数。

    \item 求幂级数 $ \sum_{n=1}^{\infty} n(x-1)^{n} $ 的收敛域及和函数。

    \item 求幂级数 $ \sum_{n=2}^{\infty}\frac{1}{(n^{2}-1)2^{n}} $ 的收敛域及和函数。

    \item 求幂级数 $ \sum_{n=1}^{\infty} (-1)^{n} \frac{n(n+1)}{2^{2n}} $ 的收敛域及和函数。
\end{enumerate}

\end{document}
